\subsection{The covering condition and finitely generated fields}

\begin{claim}{A.R2}{Proved}{Prime-to-characteristic Kummer kernels}
Let $K$ be a field finitely generated over $\mathbf F_p$, and let $G/K$
be a semiabelian variety, an extension of an abelian variety by a torus.
For the prime-to-$p$ Kummer map on rational points, the kernel is
\[
 \ker\left(G(K)\longrightarrow
       \prod_{\ell\ne p}H^1(G_K,T_\ell G)\right)=G(K)[p^\infty].
\]
For a field finitely generated over $\mathbf Q$, the analogous map
using every prime has zero kernel. No finite-generation assertion
about $K^\times$ is needed or made.
\end{claim}
\begin{proof}
For each integer $n$ invertible in $K$, the finite Kummer sequence
injects $G(K)/nG(K)$ into $H^1(G_K,G[n])$. Compatibility in $n$
identifies the displayed kernel with
$\bigcap_{(n,p)=1}nG(K)$; in characteristic zero use all positive $n$.
Passing to a finite separable extension splitting the torus preserves
any equality $x=ny$. We may therefore establish the required
implication there; the assertion that a rational point is killed by
a power of $p$ descends through the injective inclusion on points.

For a split torus, the divisor sequence on a normal proper integral
model shows that $K^\times$ modulo its constant-field units embeds in
the free abelian group on prime divisors. Its image is free abelian.
The constants are finite in positive characteristic. In characteristic
zero they form a number field; the unit theorem and the principal
fractional ideal group show again that its multiplicative group is
finite torsion plus a free abelian group. Thus $T(K)$ is free abelian
plus finite torsion. The image of $G(K)$ in the abelian quotient $A(K)$
is finitely generated. Here we use the Lang--N\'eron theorem
\cite[Theorem 2.1]{Conrad2006}: for a finitely generated regular
extension of a constant field, rational points on an abelian variety
modulo the points of its trace form a finitely generated group. The
constant-field trace has finite rational points over a finite field,
and finitely generated points over a number field.

An extension of a free abelian group plus finite torsion by a finitely
generated abelian group is itself a free abelian group plus a finitely
generated group. Indeed choose finitely many lifts of generators and
relations of the quotient. The finitely many relation values involve
only finitely many elements of a basis of the free kernel; all remaining
basis elements split off. In such a group the intersection of all
prime-to-$p$ multiples is exactly its $p$-primary torsion; the
intersection of all multiples is zero. This proves the assertion.
The converse for a $p$-primary point is immediate because every integer
prime to $p$ acts invertibly on its finite cyclic subgroup.
\end{proof}

\begin{claim}{A.R3}{Proved}{The covering condition is equivalent to collision}
Let $C/k$ be as in \statementref{A.R1}, and let $P,Q\in C(k)$.
Their sections are geometrically conjugate if and only if the
following condition holds. For every finite extension $L/k$ and every
connected finite \emph{geometrically} prime-to-$p$ Galois cover
$D\to C_L$ with a specified $P'\in D(L)$ over $P$, there is
$Q'\in D(L)$ above $Q$ such that
$[Q'-P']\in J_D(L)[p^\infty]$. When this condition holds, the lift
$Q'$ is unique. This is the condition of cote 149, scan 36, with
\emph{$p$-torsion} understood as $p$-primary torsion.
\end{claim}
\begin{proof}
Equality of section classes identifies the Galois actions on the
prime-to-$p$ fibre functor. A fixed point $P'$ in the fibre therefore
gives a fixed point $Q'$ with equal lifted section classes. For a
geometric Galois cover of order prime to $p$, the geometric open
subgroup has its full prime-to-$p$ topology: a core is obtained by
intersecting finitely many conjugate subgroups, and the resulting
quotient embeds in a product of groups of prime-to-$p$ order.
Thus the lifted conjugacy is in the correct full arithmetic quotient
of $D$. Abelianization and \statementref{A.R2} give the torsion.

We justify uniqueness independently of the theorem being proved.
Over an algebraic closure suppose two such lifts differ by a
nonidentity deck transformation $a$ of order $m$ prime to $p$.
Then $n(aR-R)$ is principal for some power $n$ of $p$. Choose $f$
with that divisor. Its norm under $\langle a\rangle$ is a constant;
scale $f$ to make the norm one. Hilbert 90 writes
$f=a(u)/u$, with a convention for the action on divisors.
Consequently $\operatorname{div}(u)-nR$ is $a$-invariant (replacing
$a$ by its inverse changes no conclusion). Every orbit of this free
action has length $m$, so the degree of an invariant divisor is
divisible by $m$. Its degree is $-n$, a contradiction.

Conversely use a cofinal system of open characteristic normal
subgroups $N$ of $\Delta_C$. They are stable under the action defined
by $s_P$, and the subgroup $N\rtimes s_P(G_k)$ defines a connected
pointed cover $D_N/k$. Its geometric deck group need not be constant.
After a finite constant extension making it constant, the assumed
condition supplies a lift $Q_N$. Uniqueness makes this lift invariant
under the original $G_k$, hence rational over $k$. Norm maps on
Jacobians send the distinguished difference class to one at every
coarser level; uniqueness gives compatibility. Compatible rational
lifts determine a $G_k$-fixed element of the profinite path torsor
from $P$ to $Q$, equivalently a geometric conjugacy of their sections.
This proves the converse in the whole group. By \statementref{A.R1}
the condition in fact forces $P=Q$.
\end{proof}

\begin{claim}{A.R4}{Proved}{Specialization from finite type to a finite field}
If proper hyperbolic curves over all finite fields of characteristic
$p$ have injective geometrically prime-to-$p$ point-section maps,
then so do proper hyperbolic curves over every field $K$ finitely
generated over $\mathbf F_p$.
\end{claim}
\begin{proof}
Suppose $P\ne Q$ collide on $C/K$. Spread the smooth proper curve,
the two sections, and its constant genus over a smooth connected
finite-type scheme $S/\mathbf F_q$ with function field $K$, shrinking
$S$ so that the sections are disjoint. Choose a closed point $s$,
let $R=\mathcal O^h_{S,s}$ and $L=\operatorname{Frac}(R)$. Pullback
along $G_L\to G_K$ preserves the proposed geometric conjugacy.

Smooth proper specialization identifies the geometric prime-to-$p$
fundamental groups of the generic and special fibres over this
henselian base. We use \cite[Expos\'e X, Corollary 3.9, and
Expos\'e XIII, Corollaries 2.9--2.10]{SGA1}, with a smooth proper
curve and no boundary or wild inertia. The comparison is equivariant
for the arithmetic map
$G_L\to\pi_1(\operatorname{Spec}R)=G_{k(s)}$ and gives a map of
arithmetic extensions. This arithmetic map is surjective: a connected
finite \emph{\'etale} cover of the normal henselian domain $R$ has
connected generic fibre, so every finite quotient of $G_{k(s)}$
is already seen by $G_L$.

The specialization map carries the sections supplied by $P,Q$ to
those supplied by $P_s,Q_s$ after this surjection. A single geometric
conjugator therefore conjugates their values for every element of
$G_{k(s)}$. Since $P_s\ne Q_s$, this contradicts the finite-field
hypothesis. This explains the direction of specialization and the
arithmetic surjectivity needed in the reduction.
\end{proof}

\begin{claim}{A.R5}{Proved}{Every hyperbolic curve over a finitely generated field}
Let $K$ be finitely generated over $\mathbf F_p$, and let $X/K$ be
a smooth geometrically connected hyperbolic curve. Its map from
$X(K)$ to geometric conjugacy classes of sections of its arithmetic
geometrically prime-to-$p$ group is injective. For $K$ finitely
generated over $\mathbf Q$, the same assertion holds for the full
geometric fundamental group. The positive-characteristic assertion
is equivalent to \statementref{A.R1}.
\end{claim}
\begin{proof}
First make a finite algebraic extension over which the smooth
geometric compactification and its distinct boundary points descend.
This is possible because their equations and morphisms use finitely
many algebraic coefficients. The purely inseparable part induces
universal homeomorphisms and equivalences on finite \emph{\'{e}tale}
covers, and identifies the absolute Galois groups. The separable part
restricts both sections along the same Galois homomorphism. Either
part therefore preserves a collision, and distinct points remain
distinct. The enlarged field is still finitely generated. This step
handles imperfect fields and boundary residue fields that cannot be
split by a separable extension alone.

For compactified genus at least two, the quotient induced by
$X\hookrightarrow\overline X$ carries a collision to one on the
proper curve. Statements \statementref{A.R1} and
\statementref{A.R4} exclude it.

If the compactification has genus zero, make a finite separable
extension splitting three boundary points, and choose their
coordinates to be $0,1,\infty$. The coordinate is a morphism
$X\to\mathbf G_m$ injective on geometric points. A collision
has equal Kummer coordinates. By \statementref{A.R2} their ratio
is $p$-primary torsion in $K^\times$, hence one in characteristic $p$.

For compactified genus one, a putative rational point $P$ makes
the compactification an elliptic curve $E$, with origin $P$ and
nonempty boundary $B$. A second colliding point $Q$ has class of
exact order $p^a$ in $E(K)$, by \statementref{A.R2}. Choose $m$
prime to $p$ with $m^2>p^a$. Pull back the punctured curve by
$[m]:E\to E$, and make a finite separable extension splitting the
boundary and deck group. The lifted collision supplies $P'=0,Q'$
in $D=E-[m]^{-1}B$. Its class in the generalized Jacobian $G_D$ is
$p$-primary. Its projection to $E$ has exact order $p^a$, because
$[m]$ is invertible on geometric $p$-primary torsion. The torus
kernel of $G_D\to E$ has no nonzero geometric $p$-primary points,
so the generalized class itself has exact order $p^a$.
The definition of the generalized Jacobian with reduced modulus
then supplies a function $f$ with
\[
 \operatorname{div}(f)=p^a(Q'-P'),\qquad
 f(T)=1\quad(T\in[m]^{-1}B).
\]
But $f-1$ has at least $m^2$ distinct zeros and pole divisor of
degree $p^a$, impossible. The boundary points are distinct because
$[m]$ is \emph{\'etale}, and none is $P'$ or $Q'$.

In characteristic zero use the Abel map into the generalized
Jacobian and the all-prime case of \statementref{A.R2}. Its
injectivity on points follows directly: equality would give a
function with divisor $Q-P$ and value one on the reduced boundary.
For positive compactified genus a degree-one function is impossible;
for genus zero hyperbolicity supplies at least three boundary points,
while a degree-one function takes a specified value only once.
Finally finite fields are finitely generated and proper genus at
least two is hyperbolic, proving the equivalence. This answers
cote 149, scans 35--36, including the affine cases.
\end{proof}

